\documentclass[11pt]{report}
\input{preambolo-verifica-dsa.tex}

\begin{document}
\chapter*{Verifica sulle equazioni di II grado incomplete}
\begin{flushright}
    \textbf{Fila C}\\
    Tempo stimato: 70 minuti.
\end{flushright}

\begin{flushright}
    Classe: 3B \quad Data: Lunedì, 28 Settembre 2026
\end{flushright} 

\vspace{0.8em}
\regoladoro

\section*{Risolvi questi esercizi per arrivare fino a 9}
\begin{enumerate}
    \item $7+2x^2=7$
    \item $3y^2-1=11$
    \item $x^2=x$
    \item $4z^2+5=1$
    \item $(x+3)^2=9+7x$
    \item $-(x-1)^2=1$
\end{enumerate}
\section*{Risolvi questo esercizio per aggiungere punti}
È un esercizio leggermente diverso da quelli visti in classe. Prova a risolverlo \ul{senza utilizzare la formula per le equazioni di secondo grado complete}.
\begin{enumerate}
    \item[7.] $-2(x^2+2)(x+2)(x^4+1)(x^2-1)=0$ 
\end{enumerate}
\end{document}
